% =========================================================
% ppt/src/minggu-11.tex
% Slide pertemuan 11
% Bab XII — Tujuh Mode Pengalamatan
%
% Sub-CPMK 3.3 | CPMK03 | CPL03
% =========================================================

\documentclass[aspectratio=169,11pt]{beamer}
% =========================================================
% ppt/src/preamble-ppt.tex
% Tema, warna, paket, dan makro bersama seluruh deck slide.
%
% Dipakai oleh setiap berkas minggu-NN.tex dengan pola:
%   \documentclass[aspectratio=169,11pt]{beamer}
%   % =========================================================
% ppt/src/preamble-ppt.tex
% Tema, warna, paket, dan makro bersama seluruh deck slide.
%
% Dipakai oleh setiap berkas minggu-NN.tex dengan pola:
%   \documentclass[aspectratio=169,11pt]{beamer}
%   % =========================================================
% ppt/src/preamble-ppt.tex
% Tema, warna, paket, dan makro bersama seluruh deck slide.
%
% Dipakai oleh setiap berkas minggu-NN.tex dengan pola:
%   \documentclass[aspectratio=169,11pt]{beamer}
%   \input{preamble-ppt}
%
% Palet warna disalin dari buku_ajar/config/styles.tex agar
% identitas visual slide sama dengan buku ajar.
% =========================================================

% ---------------------------------------------------------
% 1. Paket
% ---------------------------------------------------------
\usepackage[T1]{fontenc}
\usepackage[utf8]{inputenc}
\usepackage[indonesian]{babel}
\usepackage{lmodern}
\usepackage{tikz}
\usetikzlibrary{arrows.meta, positioning, calc, shapes.geometric}
\usepackage{booktabs}
\usepackage{array}
\usepackage{amsmath}
\usepackage{amsfonts}
\usepackage{pifont}          % \ding untuk tanda checklist
\usepackage{listings}

% ---------------------------------------------------------
% 2. Palet warna OBE (identik dengan buku)
% ---------------------------------------------------------
\definecolor{obeBlue}{HTML}{1F3864}
\definecolor{obeBlueLight}{HTML}{DCE4F5}
\definecolor{obeTeal}{HTML}{0F6E6E}
\definecolor{obeTealLight}{HTML}{D9F0EE}
\definecolor{obeGreen}{HTML}{2E6B34}
\definecolor{obeGreenLight}{HTML}{DEEFDF}
\definecolor{obeAmber}{HTML}{8A5A00}
\definecolor{obeAmberLight}{HTML}{FDF1D6}
\definecolor{obeMaroon}{HTML}{7A2233}
\definecolor{obeMaroonLight}{HTML}{F7E1E5}
\definecolor{obeGray}{HTML}{4A4A4A}
\definecolor{obeGrayLight}{HTML}{F0F0F0}
\definecolor{codeBack}{HTML}{F7F7F4}
\definecolor{codeFrame}{HTML}{C9C9C0}

% ---------------------------------------------------------
% 3. Identitas mata kuliah (dipakai seluruh deck)
% ---------------------------------------------------------
\newcommand{\NamaMK}{Bahasa Pemrograman Tingkat Rendah}
\newcommand{\NamaMKShort}{Bhs.\ Pemrograman Tingkat Rendah}
\newcommand{\KodeMK}{MWCA010}
\newcommand{\Prodi}{Program Studi Teknik Informatika}
\newcommand{\Fakultas}{Fakultas Teknologi Informasi}
\newcommand{\PerguruanTinggi}{Universitas Bale Bandung}
\newcommand{\DosenPengampu}{Cecep Suwanda, S.Si., M.Kom.}
\newcommand{\SemesterTeks}{Semester Ganjil}

% \PertemuanTeks dan \JudulBab ditetapkan di masing-masing
% berkas minggu-NN.tex sebelum \begin{document}.
\providecommand{\PertemuanTeks}{Pertemuan}

% ---------------------------------------------------------
% 4. Tema beamer
% ---------------------------------------------------------
\usetheme{default}
\setbeamertemplate{navigation symbols}{}

\setbeamercolor{normal text}{fg=black!88}
\setbeamercolor{structure}{fg=obeBlue}
\setbeamercolor{alerted text}{fg=obeMaroon}

% --- judul slide: pita biru penuh lebar ---
\setbeamercolor{frametitle}{fg=white,bg=obeBlue}
\setbeamerfont{frametitle}{series=\bfseries,size=\large}
\setbeamertemplate{frametitle}{%
    \nointerlineskip%
    \begin{beamercolorbox}[%
        wd=\paperwidth,ht=2.9ex,dp=1.5ex,
        leftskip=0.4cm,rightskip=0.4cm]{frametitle}%
        \usebeamerfont{frametitle}\insertframetitle%
    \end{beamercolorbox}%
    \nointerlineskip%
    \hbox{\textcolor{obeTeal}{\rule{\paperwidth}{1.1pt}}}%
}

% --- kaki slide: identitas mata kuliah, pertemuan, nomor ---
\setbeamercolor{footline}{fg=white,bg=obeGray}
\setbeamerfont{footline}{size=\scriptsize,family=\sffamily}
\setbeamertemplate{footline}{%
    \leavevmode%
    \hbox{%
        \begin{beamercolorbox}[wd=0.42\paperwidth,ht=2.1ex,dp=1.1ex,
            leftskip=0.4cm]{footline}%
            \usebeamerfont{footline}%
            \KodeMK\ \textbullet\ \NamaMKShort%
        \end{beamercolorbox}%
        \begin{beamercolorbox}[wd=0.32\paperwidth,ht=2.1ex,dp=1.1ex,
            center]{footline}%
            \usebeamerfont{footline}\PertemuanTeks%
        \end{beamercolorbox}%
        \begin{beamercolorbox}[wd=0.26\paperwidth,ht=2.1ex,dp=1.1ex,
            rightskip=0.4cm]{footline}%
            \usebeamerfont{footline}%
            \insertframenumber\,/\,\inserttotalframenumber%
        \end{beamercolorbox}%
    }%
    \vskip0pt%
}

% --- daftar ---
\setbeamercolor{itemize item}{fg=obeTeal}
\setbeamercolor{itemize subitem}{fg=obeTeal}
\setbeamercolor{enumerate item}{fg=obeBlue}
\setbeamertemplate{itemize item}{\raise0.2ex\hbox{\tiny$\blacksquare$}}
\setbeamertemplate{itemize subitem}{\raise0.2ex\hbox{\scriptsize$\bullet$}}

% --- kotak (pengganti tcolorbox pada buku) ---
\setbeamercolor{block title}{bg=obeBlueLight,fg=obeBlue}
\setbeamercolor{block body}{bg=obeBlueLight!45,fg=black!88}
\setbeamercolor{block title alerted}{bg=obeMaroonLight,fg=obeMaroon}
\setbeamercolor{block body alerted}{bg=obeMaroonLight!45,fg=black!88}
\setbeamercolor{block title example}{bg=obeTealLight,fg=obeTeal}
\setbeamercolor{block body example}{bg=obeTealLight!45,fg=black!88}

% ---------------------------------------------------------
% 5. Makro notasi (sama persis dengan buku)
%
% Catatan: \note TIDAK didefinisikan di sini karena beamer
% sudah memakai \note untuk catatan pembicara. Untuk catatan
% di slide dipakai \catatan.
% ---------------------------------------------------------
\newcommand{\code}[1]{\texttt{#1}}
\newcommand{\asm}[1]{\texttt{#1}}
\newcommand{\reg}[1]{\texttt{#1}}
\newcommand{\flag}[1]{\texttt{#1}}
\newcommand{\dir}[1]{\texttt{#1}}
\newcommand{\file}[1]{\texttt{#1}}
\newcommand{\eng}[1]{\textit{#1}}

\newcommand{\hex}[1]{\texttt{#1h}}
\newcommand{\bin}[1]{\texttt{#1b}}
\newcommand{\dec}[1]{\texttt{#1}}

\newcommand{\important}[1]{\textbf{Penting:} #1}
\newcommand{\catatan}[1]{\textbf{Catatan:} #1}
\newcommand{\term}[1]{\textbf{#1}}

\newcommand{\CPL}[1]{\textbf{CPL#1}}
\newcommand{\CPMK}[1]{\textbf{CPMK#1}}
\newcommand{\SubCPMK}[1]{\textbf{Sub-CPMK #1}}

\newcommand{\checkbox}{\ding{113}\hspace{0.35em}}
\newcommand{\checked}{\ding{51}\hspace{0.35em}}
\newcommand{\crossed}{\ding{55}\hspace{0.35em}}

\newcommand{\flowarrow}{\ensuremath{\rightarrow}}

% ---------------------------------------------------------
% 6. Makro tata letak slide
% ---------------------------------------------------------

% Keterangan sumber materi (ketertelusuran ke buku)
\newcommand{\Sumber}[1]{%
    \vfill
    {\scriptsize\color{obeGray}Sumber: #1\par}%
}

% Frame judul deck: \JudulDeck{<judul bab>}{<pertemuan>}
\newcommand{\JudulDeck}[2]{%
\begin{frame}[plain]
    \vspace{0.25cm}
    {\color{obeBlue}\sffamily\bfseries\small \NamaMK\ (\KodeMK)\par}
    \vspace{0.12cm}
    {\color{obeGray}\sffamily\scriptsize
     \Prodi\ \textbullet\ \Fakultas\ \textbullet\ \PerguruanTinggi\par}
    \vfill
    {\color{obeBlue}\sffamily\LARGE\bfseries #1\par}
    \vspace{0.4cm}
    {\color{obeTeal}\sffamily\large #2\par}
    \vspace{0.45cm}
    {\color{obeTeal}\rule{\textwidth}{1pt}}
    \vfill
    {\color{obeGray}\sffamily\small Dosen Pengampu\par}
    \vspace{0.08cm}
    {\sffamily\bfseries \DosenPengampu\par}
    \vspace{0.2cm}
    {\color{obeGray}\sffamily\small \SemesterTeks\par}
    \vspace{0.25cm}
\end{frame}
}

% Frame pemisah bagian
\newcommand{\PemisahBagian}[1]{%
\begin{frame}[plain]
    \vfill
    \begin{center}
        {\color{obeTeal}\sffamily\small\bfseries BAGIAN\par}
        \vspace{0.35cm}
        {\color{obeBlue}\sffamily\Large\bfseries #1\par}
        \vspace{0.35cm}
        {\color{obeTeal}\rule{0.55\textwidth}{0.9pt}}
    \end{center}
    \vfill
\end{frame}
}

% ---------------------------------------------------------
% 7. Gaya penampil kode
% ---------------------------------------------------------
\input{listing-ppt}

%
% Palet warna disalin dari buku_ajar/config/styles.tex agar
% identitas visual slide sama dengan buku ajar.
% =========================================================

% ---------------------------------------------------------
% 1. Paket
% ---------------------------------------------------------
\usepackage[T1]{fontenc}
\usepackage[utf8]{inputenc}
\usepackage[indonesian]{babel}
\usepackage{lmodern}
\usepackage{tikz}
\usetikzlibrary{arrows.meta, positioning, calc, shapes.geometric}
\usepackage{booktabs}
\usepackage{array}
\usepackage{amsmath}
\usepackage{amsfonts}
\usepackage{pifont}          % \ding untuk tanda checklist
\usepackage{listings}

% ---------------------------------------------------------
% 2. Palet warna OBE (identik dengan buku)
% ---------------------------------------------------------
\definecolor{obeBlue}{HTML}{1F3864}
\definecolor{obeBlueLight}{HTML}{DCE4F5}
\definecolor{obeTeal}{HTML}{0F6E6E}
\definecolor{obeTealLight}{HTML}{D9F0EE}
\definecolor{obeGreen}{HTML}{2E6B34}
\definecolor{obeGreenLight}{HTML}{DEEFDF}
\definecolor{obeAmber}{HTML}{8A5A00}
\definecolor{obeAmberLight}{HTML}{FDF1D6}
\definecolor{obeMaroon}{HTML}{7A2233}
\definecolor{obeMaroonLight}{HTML}{F7E1E5}
\definecolor{obeGray}{HTML}{4A4A4A}
\definecolor{obeGrayLight}{HTML}{F0F0F0}
\definecolor{codeBack}{HTML}{F7F7F4}
\definecolor{codeFrame}{HTML}{C9C9C0}

% ---------------------------------------------------------
% 3. Identitas mata kuliah (dipakai seluruh deck)
% ---------------------------------------------------------
\newcommand{\NamaMK}{Bahasa Pemrograman Tingkat Rendah}
\newcommand{\NamaMKShort}{Bhs.\ Pemrograman Tingkat Rendah}
\newcommand{\KodeMK}{MWCA010}
\newcommand{\Prodi}{Program Studi Teknik Informatika}
\newcommand{\Fakultas}{Fakultas Teknologi Informasi}
\newcommand{\PerguruanTinggi}{Universitas Bale Bandung}
\newcommand{\DosenPengampu}{Cecep Suwanda, S.Si., M.Kom.}
\newcommand{\SemesterTeks}{Semester Ganjil}

% \PertemuanTeks dan \JudulBab ditetapkan di masing-masing
% berkas minggu-NN.tex sebelum \begin{document}.
\providecommand{\PertemuanTeks}{Pertemuan}

% ---------------------------------------------------------
% 4. Tema beamer
% ---------------------------------------------------------
\usetheme{default}
\setbeamertemplate{navigation symbols}{}

\setbeamercolor{normal text}{fg=black!88}
\setbeamercolor{structure}{fg=obeBlue}
\setbeamercolor{alerted text}{fg=obeMaroon}

% --- judul slide: pita biru penuh lebar ---
\setbeamercolor{frametitle}{fg=white,bg=obeBlue}
\setbeamerfont{frametitle}{series=\bfseries,size=\large}
\setbeamertemplate{frametitle}{%
    \nointerlineskip%
    \begin{beamercolorbox}[%
        wd=\paperwidth,ht=2.9ex,dp=1.5ex,
        leftskip=0.4cm,rightskip=0.4cm]{frametitle}%
        \usebeamerfont{frametitle}\insertframetitle%
    \end{beamercolorbox}%
    \nointerlineskip%
    \hbox{\textcolor{obeTeal}{\rule{\paperwidth}{1.1pt}}}%
}

% --- kaki slide: identitas mata kuliah, pertemuan, nomor ---
\setbeamercolor{footline}{fg=white,bg=obeGray}
\setbeamerfont{footline}{size=\scriptsize,family=\sffamily}
\setbeamertemplate{footline}{%
    \leavevmode%
    \hbox{%
        \begin{beamercolorbox}[wd=0.42\paperwidth,ht=2.1ex,dp=1.1ex,
            leftskip=0.4cm]{footline}%
            \usebeamerfont{footline}%
            \KodeMK\ \textbullet\ \NamaMKShort%
        \end{beamercolorbox}%
        \begin{beamercolorbox}[wd=0.32\paperwidth,ht=2.1ex,dp=1.1ex,
            center]{footline}%
            \usebeamerfont{footline}\PertemuanTeks%
        \end{beamercolorbox}%
        \begin{beamercolorbox}[wd=0.26\paperwidth,ht=2.1ex,dp=1.1ex,
            rightskip=0.4cm]{footline}%
            \usebeamerfont{footline}%
            \insertframenumber\,/\,\inserttotalframenumber%
        \end{beamercolorbox}%
    }%
    \vskip0pt%
}

% --- daftar ---
\setbeamercolor{itemize item}{fg=obeTeal}
\setbeamercolor{itemize subitem}{fg=obeTeal}
\setbeamercolor{enumerate item}{fg=obeBlue}
\setbeamertemplate{itemize item}{\raise0.2ex\hbox{\tiny$\blacksquare$}}
\setbeamertemplate{itemize subitem}{\raise0.2ex\hbox{\scriptsize$\bullet$}}

% --- kotak (pengganti tcolorbox pada buku) ---
\setbeamercolor{block title}{bg=obeBlueLight,fg=obeBlue}
\setbeamercolor{block body}{bg=obeBlueLight!45,fg=black!88}
\setbeamercolor{block title alerted}{bg=obeMaroonLight,fg=obeMaroon}
\setbeamercolor{block body alerted}{bg=obeMaroonLight!45,fg=black!88}
\setbeamercolor{block title example}{bg=obeTealLight,fg=obeTeal}
\setbeamercolor{block body example}{bg=obeTealLight!45,fg=black!88}

% ---------------------------------------------------------
% 5. Makro notasi (sama persis dengan buku)
%
% Catatan: \note TIDAK didefinisikan di sini karena beamer
% sudah memakai \note untuk catatan pembicara. Untuk catatan
% di slide dipakai \catatan.
% ---------------------------------------------------------
\newcommand{\code}[1]{\texttt{#1}}
\newcommand{\asm}[1]{\texttt{#1}}
\newcommand{\reg}[1]{\texttt{#1}}
\newcommand{\flag}[1]{\texttt{#1}}
\newcommand{\dir}[1]{\texttt{#1}}
\newcommand{\file}[1]{\texttt{#1}}
\newcommand{\eng}[1]{\textit{#1}}

\newcommand{\hex}[1]{\texttt{#1h}}
\newcommand{\bin}[1]{\texttt{#1b}}
\newcommand{\dec}[1]{\texttt{#1}}

\newcommand{\important}[1]{\textbf{Penting:} #1}
\newcommand{\catatan}[1]{\textbf{Catatan:} #1}
\newcommand{\term}[1]{\textbf{#1}}

\newcommand{\CPL}[1]{\textbf{CPL#1}}
\newcommand{\CPMK}[1]{\textbf{CPMK#1}}
\newcommand{\SubCPMK}[1]{\textbf{Sub-CPMK #1}}

\newcommand{\checkbox}{\ding{113}\hspace{0.35em}}
\newcommand{\checked}{\ding{51}\hspace{0.35em}}
\newcommand{\crossed}{\ding{55}\hspace{0.35em}}

\newcommand{\flowarrow}{\ensuremath{\rightarrow}}

% ---------------------------------------------------------
% 6. Makro tata letak slide
% ---------------------------------------------------------

% Keterangan sumber materi (ketertelusuran ke buku)
\newcommand{\Sumber}[1]{%
    \vfill
    {\scriptsize\color{obeGray}Sumber: #1\par}%
}

% Frame judul deck: \JudulDeck{<judul bab>}{<pertemuan>}
\newcommand{\JudulDeck}[2]{%
\begin{frame}[plain]
    \vspace{0.25cm}
    {\color{obeBlue}\sffamily\bfseries\small \NamaMK\ (\KodeMK)\par}
    \vspace{0.12cm}
    {\color{obeGray}\sffamily\scriptsize
     \Prodi\ \textbullet\ \Fakultas\ \textbullet\ \PerguruanTinggi\par}
    \vfill
    {\color{obeBlue}\sffamily\LARGE\bfseries #1\par}
    \vspace{0.4cm}
    {\color{obeTeal}\sffamily\large #2\par}
    \vspace{0.45cm}
    {\color{obeTeal}\rule{\textwidth}{1pt}}
    \vfill
    {\color{obeGray}\sffamily\small Dosen Pengampu\par}
    \vspace{0.08cm}
    {\sffamily\bfseries \DosenPengampu\par}
    \vspace{0.2cm}
    {\color{obeGray}\sffamily\small \SemesterTeks\par}
    \vspace{0.25cm}
\end{frame}
}

% Frame pemisah bagian
\newcommand{\PemisahBagian}[1]{%
\begin{frame}[plain]
    \vfill
    \begin{center}
        {\color{obeTeal}\sffamily\small\bfseries BAGIAN\par}
        \vspace{0.35cm}
        {\color{obeBlue}\sffamily\Large\bfseries #1\par}
        \vspace{0.35cm}
        {\color{obeTeal}\rule{0.55\textwidth}{0.9pt}}
    \end{center}
    \vfill
\end{frame}
}

% ---------------------------------------------------------
% 7. Gaya penampil kode
% ---------------------------------------------------------
% =========================================================
% ppt/src/listing-ppt.tex
% Gaya penampil kode untuk slide.
%
% Diturunkan dari buku_ajar/config/listings.tex, tetapi
% disesuaikan untuk slide: ukuran lebih kecil (karena frame
% 16:9 dan kode harus utuh terlihat) dan jarak lebih rapat.
% =========================================================

% ---------------------------------------------------------
% 1. Definisi bahasa MASM/TASM 8086 (salinan dari buku)
% ---------------------------------------------------------
\lstdefinelanguage{MASM}{
    morekeywords={
        MOV, XCHG, LEA, LDS, LES, PUSH, POP, PUSHA, POPA, PUSHF, POPF,
        LAHF, SAHF, XLAT, IN, OUT, CBW, CWD,
        ADD, ADC, SUB, SBB, INC, DEC, NEG, CMP, MUL, IMUL, DIV, IDIV,
        AAA, AAS, AAM, AAD, DAA, DAS,
        AND, OR, XOR, NOT, TEST, SHL, SHR, SAL, SAR, ROL, ROR, RCL, RCR,
        JMP, JE, JZ, JNE, JNZ, JG, JGE, JL, JLE, JA, JAE, JB, JBE,
        JC, JNC, JS, JNS, JO, JNO, JP, JPE, JNP, JPO, JCXZ,
        LOOP, LOOPE, LOOPZ, LOOPNE, LOOPNZ, CALL, RET, INT, IRET, INTO,
        HLT, NOP, CLI, STI, CLD, STD, CLC, STC, CMC, WAIT, ESC, LOCK,
        DB, DW, DD, DQ, DT, EQU, ORG, SEGMENT, ENDS, ASSUME, END,
        PROC, ENDP, MACRO, ENDM, LOCAL, PUBLIC, EXTRN, EXTERN,
        INCLUDE, INCLUDELIB, MODEL, STACK, CODE, DATA, CONST, ALIGN, EVEN,
        OFFSET, PTR, BYTE, WORD, DWORD, QWORD, TBYTE,
        NEAR, FAR, SHORT, LABEL, TYPE, SIZE, LENGTH, DUP, THIS,
        AX, BX, CX, DX, AL, AH, BL, BH, CL, CH, DL, DH,
        SP, BP, SI, DI, IP, CS, DS, SS, ES, FLAGS,
        CF, PF, AF, ZF, TF, IF, DF, OF
    },
    sensitive=false,
    morecomment=[l]{;},
    morestring=[b]',
    morestring=[b]"
}

% ---------------------------------------------------------
% 2. Gaya kode assembly untuk slide
% ---------------------------------------------------------
\lstdefinestyle{asmslide}{
    language=MASM,
    basicstyle=\ttfamily\scriptsize,
    keywordstyle=\color{obeBlue}\bfseries,
    commentstyle=\color{obeGreen}\itshape,
    stringstyle=\color{obeMaroon},
    numberstyle=\tiny\color{obeGray},
    numbers=left,
    numbersep=5pt,
    stepnumber=1,
    backgroundcolor=\color{codeBack},
    frame=single,
    rulecolor=\color{codeFrame},
    framesep=4pt,
    breaklines=true,
    breakatwhitespace=false,
    % CATATAN PENTING: opsi `postbreak` TIDAK boleh dipakai di sini.
    % Di beamer, `breaklines=true` bersama `postbreak=<kode>` membuat
    % listings gagal dirakit dengan galat
    %   "Missing number, treated as zero <to be read again> \kern".
    % Karena \lstset{style=asmslide} bersifat global, opsi postbreak
    % juga bocor ke gaya lain (mis. outslide/keluaran) dan merusaknya.
    % Di buku ajar (kelas `book`) justru sebaliknya: postbreak aman dan
    % dipakai. Jadi larangan ini khusus untuk slide.
    showstringspaces=false,
    tabsize=4,
    columns=fullflexible,
    keepspaces=true,
    upquote=true,
    xleftmargin=12pt,
    aboveskip=0.4em,
    belowskip=0.4em
}

% ---------------------------------------------------------
% 3. Gaya keluaran program / sesi debugger untuk slide
% ---------------------------------------------------------
\lstdefinestyle{outslide}{
    basicstyle=\ttfamily\scriptsize,
    keywordstyle=\color{black},
    commentstyle=\color{obeGray},
    stringstyle=\color{black},
    numbers=none,
    backgroundcolor=\color{obeGrayLight},
    frame=single,
    rulecolor=\color{codeFrame},
    framesep=4pt,
    breaklines=true,
    showstringspaces=false,
    columns=fullflexible,
    keepspaces=true,
    upquote=true,
    xleftmargin=4pt,
    aboveskip=0.4em,
    belowskip=0.4em
}

\lstset{style=asmslide}

% ---------------------------------------------------------
% 4. Environment ringkas untuk dipakai di dalam frame
%
% Frame yang memuat kedua environment ini WAJIB memakai
% opsi [fragile], karena isinya verbatim.
% ---------------------------------------------------------
\lstnewenvironment{kodeasm}[1][]{%
    \lstset{style=asmslide, #1}%
}{}

\lstnewenvironment{keluaran}[1][]{%
    \lstset{style=outslide, #1}%
}{}


%
% Palet warna disalin dari buku_ajar/config/styles.tex agar
% identitas visual slide sama dengan buku ajar.
% =========================================================

% ---------------------------------------------------------
% 1. Paket
% ---------------------------------------------------------
\usepackage[T1]{fontenc}
\usepackage[utf8]{inputenc}
\usepackage[indonesian]{babel}
\usepackage{lmodern}
\usepackage{tikz}
\usetikzlibrary{arrows.meta, positioning, calc, shapes.geometric}
\usepackage{booktabs}
\usepackage{array}
\usepackage{amsmath}
\usepackage{amsfonts}
\usepackage{pifont}          % \ding untuk tanda checklist
\usepackage{listings}

% ---------------------------------------------------------
% 2. Palet warna OBE (identik dengan buku)
% ---------------------------------------------------------
\definecolor{obeBlue}{HTML}{1F3864}
\definecolor{obeBlueLight}{HTML}{DCE4F5}
\definecolor{obeTeal}{HTML}{0F6E6E}
\definecolor{obeTealLight}{HTML}{D9F0EE}
\definecolor{obeGreen}{HTML}{2E6B34}
\definecolor{obeGreenLight}{HTML}{DEEFDF}
\definecolor{obeAmber}{HTML}{8A5A00}
\definecolor{obeAmberLight}{HTML}{FDF1D6}
\definecolor{obeMaroon}{HTML}{7A2233}
\definecolor{obeMaroonLight}{HTML}{F7E1E5}
\definecolor{obeGray}{HTML}{4A4A4A}
\definecolor{obeGrayLight}{HTML}{F0F0F0}
\definecolor{codeBack}{HTML}{F7F7F4}
\definecolor{codeFrame}{HTML}{C9C9C0}

% ---------------------------------------------------------
% 3. Identitas mata kuliah (dipakai seluruh deck)
% ---------------------------------------------------------
\newcommand{\NamaMK}{Bahasa Pemrograman Tingkat Rendah}
\newcommand{\NamaMKShort}{Bhs.\ Pemrograman Tingkat Rendah}
\newcommand{\KodeMK}{MWCA010}
\newcommand{\Prodi}{Program Studi Teknik Informatika}
\newcommand{\Fakultas}{Fakultas Teknologi Informasi}
\newcommand{\PerguruanTinggi}{Universitas Bale Bandung}
\newcommand{\DosenPengampu}{Cecep Suwanda, S.Si., M.Kom.}
\newcommand{\SemesterTeks}{Semester Ganjil}

% \PertemuanTeks dan \JudulBab ditetapkan di masing-masing
% berkas minggu-NN.tex sebelum \begin{document}.
\providecommand{\PertemuanTeks}{Pertemuan}

% ---------------------------------------------------------
% 4. Tema beamer
% ---------------------------------------------------------
\usetheme{default}
\setbeamertemplate{navigation symbols}{}

\setbeamercolor{normal text}{fg=black!88}
\setbeamercolor{structure}{fg=obeBlue}
\setbeamercolor{alerted text}{fg=obeMaroon}

% --- judul slide: pita biru penuh lebar ---
\setbeamercolor{frametitle}{fg=white,bg=obeBlue}
\setbeamerfont{frametitle}{series=\bfseries,size=\large}
\setbeamertemplate{frametitle}{%
    \nointerlineskip%
    \begin{beamercolorbox}[%
        wd=\paperwidth,ht=2.9ex,dp=1.5ex,
        leftskip=0.4cm,rightskip=0.4cm]{frametitle}%
        \usebeamerfont{frametitle}\insertframetitle%
    \end{beamercolorbox}%
    \nointerlineskip%
    \hbox{\textcolor{obeTeal}{\rule{\paperwidth}{1.1pt}}}%
}

% --- kaki slide: identitas mata kuliah, pertemuan, nomor ---
\setbeamercolor{footline}{fg=white,bg=obeGray}
\setbeamerfont{footline}{size=\scriptsize,family=\sffamily}
\setbeamertemplate{footline}{%
    \leavevmode%
    \hbox{%
        \begin{beamercolorbox}[wd=0.42\paperwidth,ht=2.1ex,dp=1.1ex,
            leftskip=0.4cm]{footline}%
            \usebeamerfont{footline}%
            \KodeMK\ \textbullet\ \NamaMKShort%
        \end{beamercolorbox}%
        \begin{beamercolorbox}[wd=0.32\paperwidth,ht=2.1ex,dp=1.1ex,
            center]{footline}%
            \usebeamerfont{footline}\PertemuanTeks%
        \end{beamercolorbox}%
        \begin{beamercolorbox}[wd=0.26\paperwidth,ht=2.1ex,dp=1.1ex,
            rightskip=0.4cm]{footline}%
            \usebeamerfont{footline}%
            \insertframenumber\,/\,\inserttotalframenumber%
        \end{beamercolorbox}%
    }%
    \vskip0pt%
}

% --- daftar ---
\setbeamercolor{itemize item}{fg=obeTeal}
\setbeamercolor{itemize subitem}{fg=obeTeal}
\setbeamercolor{enumerate item}{fg=obeBlue}
\setbeamertemplate{itemize item}{\raise0.2ex\hbox{\tiny$\blacksquare$}}
\setbeamertemplate{itemize subitem}{\raise0.2ex\hbox{\scriptsize$\bullet$}}

% --- kotak (pengganti tcolorbox pada buku) ---
\setbeamercolor{block title}{bg=obeBlueLight,fg=obeBlue}
\setbeamercolor{block body}{bg=obeBlueLight!45,fg=black!88}
\setbeamercolor{block title alerted}{bg=obeMaroonLight,fg=obeMaroon}
\setbeamercolor{block body alerted}{bg=obeMaroonLight!45,fg=black!88}
\setbeamercolor{block title example}{bg=obeTealLight,fg=obeTeal}
\setbeamercolor{block body example}{bg=obeTealLight!45,fg=black!88}

% ---------------------------------------------------------
% 5. Makro notasi (sama persis dengan buku)
%
% Catatan: \note TIDAK didefinisikan di sini karena beamer
% sudah memakai \note untuk catatan pembicara. Untuk catatan
% di slide dipakai \catatan.
% ---------------------------------------------------------
\newcommand{\code}[1]{\texttt{#1}}
\newcommand{\asm}[1]{\texttt{#1}}
\newcommand{\reg}[1]{\texttt{#1}}
\newcommand{\flag}[1]{\texttt{#1}}
\newcommand{\dir}[1]{\texttt{#1}}
\newcommand{\file}[1]{\texttt{#1}}
\newcommand{\eng}[1]{\textit{#1}}

\newcommand{\hex}[1]{\texttt{#1h}}
\newcommand{\bin}[1]{\texttt{#1b}}
\newcommand{\dec}[1]{\texttt{#1}}

\newcommand{\important}[1]{\textbf{Penting:} #1}
\newcommand{\catatan}[1]{\textbf{Catatan:} #1}
\newcommand{\term}[1]{\textbf{#1}}

\newcommand{\CPL}[1]{\textbf{CPL#1}}
\newcommand{\CPMK}[1]{\textbf{CPMK#1}}
\newcommand{\SubCPMK}[1]{\textbf{Sub-CPMK #1}}

\newcommand{\checkbox}{\ding{113}\hspace{0.35em}}
\newcommand{\checked}{\ding{51}\hspace{0.35em}}
\newcommand{\crossed}{\ding{55}\hspace{0.35em}}

\newcommand{\flowarrow}{\ensuremath{\rightarrow}}

% ---------------------------------------------------------
% 6. Makro tata letak slide
% ---------------------------------------------------------

% Keterangan sumber materi (ketertelusuran ke buku)
\newcommand{\Sumber}[1]{%
    \vfill
    {\scriptsize\color{obeGray}Sumber: #1\par}%
}

% Frame judul deck: \JudulDeck{<judul bab>}{<pertemuan>}
\newcommand{\JudulDeck}[2]{%
\begin{frame}[plain]
    \vspace{0.25cm}
    {\color{obeBlue}\sffamily\bfseries\small \NamaMK\ (\KodeMK)\par}
    \vspace{0.12cm}
    {\color{obeGray}\sffamily\scriptsize
     \Prodi\ \textbullet\ \Fakultas\ \textbullet\ \PerguruanTinggi\par}
    \vfill
    {\color{obeBlue}\sffamily\LARGE\bfseries #1\par}
    \vspace{0.4cm}
    {\color{obeTeal}\sffamily\large #2\par}
    \vspace{0.45cm}
    {\color{obeTeal}\rule{\textwidth}{1pt}}
    \vfill
    {\color{obeGray}\sffamily\small Dosen Pengampu\par}
    \vspace{0.08cm}
    {\sffamily\bfseries \DosenPengampu\par}
    \vspace{0.2cm}
    {\color{obeGray}\sffamily\small \SemesterTeks\par}
    \vspace{0.25cm}
\end{frame}
}

% Frame pemisah bagian
\newcommand{\PemisahBagian}[1]{%
\begin{frame}[plain]
    \vfill
    \begin{center}
        {\color{obeTeal}\sffamily\small\bfseries BAGIAN\par}
        \vspace{0.35cm}
        {\color{obeBlue}\sffamily\Large\bfseries #1\par}
        \vspace{0.35cm}
        {\color{obeTeal}\rule{0.55\textwidth}{0.9pt}}
    \end{center}
    \vfill
\end{frame}
}

% ---------------------------------------------------------
% 7. Gaya penampil kode
% ---------------------------------------------------------
% =========================================================
% ppt/src/listing-ppt.tex
% Gaya penampil kode untuk slide.
%
% Diturunkan dari buku_ajar/config/listings.tex, tetapi
% disesuaikan untuk slide: ukuran lebih kecil (karena frame
% 16:9 dan kode harus utuh terlihat) dan jarak lebih rapat.
% =========================================================

% ---------------------------------------------------------
% 1. Definisi bahasa MASM/TASM 8086 (salinan dari buku)
% ---------------------------------------------------------
\lstdefinelanguage{MASM}{
    morekeywords={
        MOV, XCHG, LEA, LDS, LES, PUSH, POP, PUSHA, POPA, PUSHF, POPF,
        LAHF, SAHF, XLAT, IN, OUT, CBW, CWD,
        ADD, ADC, SUB, SBB, INC, DEC, NEG, CMP, MUL, IMUL, DIV, IDIV,
        AAA, AAS, AAM, AAD, DAA, DAS,
        AND, OR, XOR, NOT, TEST, SHL, SHR, SAL, SAR, ROL, ROR, RCL, RCR,
        JMP, JE, JZ, JNE, JNZ, JG, JGE, JL, JLE, JA, JAE, JB, JBE,
        JC, JNC, JS, JNS, JO, JNO, JP, JPE, JNP, JPO, JCXZ,
        LOOP, LOOPE, LOOPZ, LOOPNE, LOOPNZ, CALL, RET, INT, IRET, INTO,
        HLT, NOP, CLI, STI, CLD, STD, CLC, STC, CMC, WAIT, ESC, LOCK,
        DB, DW, DD, DQ, DT, EQU, ORG, SEGMENT, ENDS, ASSUME, END,
        PROC, ENDP, MACRO, ENDM, LOCAL, PUBLIC, EXTRN, EXTERN,
        INCLUDE, INCLUDELIB, MODEL, STACK, CODE, DATA, CONST, ALIGN, EVEN,
        OFFSET, PTR, BYTE, WORD, DWORD, QWORD, TBYTE,
        NEAR, FAR, SHORT, LABEL, TYPE, SIZE, LENGTH, DUP, THIS,
        AX, BX, CX, DX, AL, AH, BL, BH, CL, CH, DL, DH,
        SP, BP, SI, DI, IP, CS, DS, SS, ES, FLAGS,
        CF, PF, AF, ZF, TF, IF, DF, OF
    },
    sensitive=false,
    morecomment=[l]{;},
    morestring=[b]',
    morestring=[b]"
}

% ---------------------------------------------------------
% 2. Gaya kode assembly untuk slide
% ---------------------------------------------------------
\lstdefinestyle{asmslide}{
    language=MASM,
    basicstyle=\ttfamily\scriptsize,
    keywordstyle=\color{obeBlue}\bfseries,
    commentstyle=\color{obeGreen}\itshape,
    stringstyle=\color{obeMaroon},
    numberstyle=\tiny\color{obeGray},
    numbers=left,
    numbersep=5pt,
    stepnumber=1,
    backgroundcolor=\color{codeBack},
    frame=single,
    rulecolor=\color{codeFrame},
    framesep=4pt,
    breaklines=true,
    breakatwhitespace=false,
    % CATATAN PENTING: opsi `postbreak` TIDAK boleh dipakai di sini.
    % Di beamer, `breaklines=true` bersama `postbreak=<kode>` membuat
    % listings gagal dirakit dengan galat
    %   "Missing number, treated as zero <to be read again> \kern".
    % Karena \lstset{style=asmslide} bersifat global, opsi postbreak
    % juga bocor ke gaya lain (mis. outslide/keluaran) dan merusaknya.
    % Di buku ajar (kelas `book`) justru sebaliknya: postbreak aman dan
    % dipakai. Jadi larangan ini khusus untuk slide.
    showstringspaces=false,
    tabsize=4,
    columns=fullflexible,
    keepspaces=true,
    upquote=true,
    xleftmargin=12pt,
    aboveskip=0.4em,
    belowskip=0.4em
}

% ---------------------------------------------------------
% 3. Gaya keluaran program / sesi debugger untuk slide
% ---------------------------------------------------------
\lstdefinestyle{outslide}{
    basicstyle=\ttfamily\scriptsize,
    keywordstyle=\color{black},
    commentstyle=\color{obeGray},
    stringstyle=\color{black},
    numbers=none,
    backgroundcolor=\color{obeGrayLight},
    frame=single,
    rulecolor=\color{codeFrame},
    framesep=4pt,
    breaklines=true,
    showstringspaces=false,
    columns=fullflexible,
    keepspaces=true,
    upquote=true,
    xleftmargin=4pt,
    aboveskip=0.4em,
    belowskip=0.4em
}

\lstset{style=asmslide}

% ---------------------------------------------------------
% 4. Environment ringkas untuk dipakai di dalam frame
%
% Frame yang memuat kedua environment ini WAJIB memakai
% opsi [fragile], karena isinya verbatim.
% ---------------------------------------------------------
\lstnewenvironment{kodeasm}[1][]{%
    \lstset{style=asmslide, #1}%
}{}

\lstnewenvironment{keluaran}[1][]{%
    \lstset{style=outslide, #1}%
}{}



\renewcommand{\PertemuanTeks}{Pertemuan 11}

\begin{document}

% =========================================================
% 1. JUDUL
% =========================================================
\JudulDeck{Tujuh Mode Pengalamatan}{Pertemuan 11}

% =========================================================
% 2. CAPAIAN DAN TUJUAN
% =========================================================
\begin{frame}{Capaian dan Tujuan Pembelajaran}
\begin{block}{Sub-CPMK 3.3}
\footnotesize Menerapkan tujuh mode pengalamatan 8086
(\eng{register}, \eng{immediate}, \eng{direct memory},
\eng{indirect memory}, \eng{base relative}, \eng{direct index},
\eng{based indexed}) pada instruksi transfer dan aritmatika.
\end{block}

\vspace{0.1em}
{\small\textbf{CPMK yang didukung:} \CPMK{03}\quad\textbullet\quad
\textbf{CPL yang didukung:} \CPL{03}\par}

\vspace{0.15em}
\begin{exampleblock}{Setelah pertemuan ini, mahasiswa mampu:}
\footnotesize
\begin{enumerate}
    \item menjelaskan peran mode pengalamatan;
    \item mengenali mode dari penulisan operand;
    \item menerapkan mode register, immediate, dan direct memory;
    \item menerapkan indirect, base relative, direct index, dan
          based indexed;
    \item menghitung alamat efektif/fisik serta memilih mode untuk
          array dan \eng{record}.
\end{enumerate}
\end{exampleblock}
\end{frame}

% =========================================================
% 3. PETA MATERI
% =========================================================
\begin{frame}{Peta Materi}
\small
\begin{enumerate}
    \item \textbf{Operand dan Tiga Mode Dasar}\\
          {\footnotesize operand sumber/tujuan; register; immediate;
          direct memory}

    \vspace{0.25em}
    \item \textbf{Mode Memori Berpenunjuk}\\
          {\footnotesize indirect; base relative; direct index;
          based indexed; alamat efektif; kombinasi legal}

    \vspace{0.25em}
    \item \textbf{Segmen Default dan Penerapan}\\
          {\footnotesize segmen default; penggantian segmen; memilih
          mode untuk array dan record}
\end{enumerate}

\vspace{0.35em}
\begin{block}{Praktik pada pertemuan ini}
\small Membandingkan mode pengalamatan pada akses array yang sama.
\end{block}
\end{frame}

% =========================================================
% BAGIAN 1
% =========================================================
\PemisahBagian{Operand dan Tiga Mode Dasar}

% ---------------------------------------------------------
\begin{frame}{Operand Sumber dan Tujuan}
\small
\begin{itemize}
    \item Sebagian besar instruksi 8086 memakai dua operand: sumber
          dan tujuan.

    \vspace{0.18em}
    \item Tujuan harus berupa register atau lokasi memori; nilai
          segera tidak dapat menjadi tujuan.

    \vspace{0.18em}
    \item Mode pengalamatan menyatakan di mana operand berada dan
          bagaimana prosesor memperoleh nilainya.

    \vspace{0.18em}
    \item Contoh: \code{MOV AX, BX} dan \code{MOV AX, 1234h}.
\end{itemize}

\vspace{0.25em}
\begin{block}{Tiga tempat operand}
\small Di dalam instruksi, di dalam register, atau di dalam memori.
\end{block}
\Sumber{Bab XII, Bagian 1}
\end{frame}

% ---------------------------------------------------------
\begin{frame}{Tujuh Mode Pengalamatan}
\scriptsize
\begin{center}
\begin{tabular}{@{}>{\raggedright\arraybackslash}p{2.4cm}
    >{\raggedright\arraybackslash}p{3.6cm}
    >{\raggedright\arraybackslash}p{3.4cm}
    >{\raggedright\arraybackslash}p{3.0cm}@{}}
    \toprule
    \textbf{Mode} & \textbf{Bentuk penulisan} & \textbf{Contoh} &
    \textbf{Operand di} \\
    \midrule
    Register & register & \code{MOV AX, BX} & register \\
    \addlinespace
    Immediate & nilai segera & \code{MOV AX, 1234h} &
    dalam instruksi \\
    \addlinespace
    Direct & \code{[alamat]} atau nama & \code{MOV AX, [1234h]} &
    memori, alamat tetap \\
    \addlinespace
    Indirect & \code{[register]} & \code{MOV AX, [BX]} &
    memori, alamat dari register \\
    \addlinespace
    Base relative & \code{[basis+disp]} & \code{MOV AX, [BX+4]} &
    memori, basis + disp \\
    \addlinespace
    Direct index & \code{[indeks+disp]} & \code{MOV AX, [SI+4]} &
    memori, indeks + disp \\
    \addlinespace
    Based indexed & \code{[basis+indeks+disp]} &
    \code{MOV AX, [BX+SI+4]} & memori, basis + indeks + disp \\
    \bottomrule
\end{tabular}
\end{center}
\Sumber{Bab XII, Bagian 1}
\end{frame}

% ---------------------------------------------------------
\begin{frame}[fragile]{Mode Register}
\small
Operand berada di dalam register; tidak ada pengaksesan memori.

\vspace{0.2em}
\begin{kodeasm}[numbers=none]
MOV AX, BX          ; AX <- BX
ADD DX, CX          ; DX <- DX + CX
MOV DL, AL          ; DL <- AL
\end{kodeasm}

\vspace{0.15em}
{\footnotesize\color{obeGray}Paling cepat dan ringkas, tetapi jumlah
register 8086 terbatas.\par}
\Sumber{Bab XII, Bagian 1}
\end{frame}

% ---------------------------------------------------------
\begin{frame}[fragile]{Mode Immediate}
\small
Operand sumber adalah nilai konstan yang tertulis di dalam instruksi;
tidak perlu mengakses memori.

\vspace{0.2em}
\begin{kodeasm}[numbers=none]
MOV AX, 1234h       ; immediate 16 bit
MOV BL, 0Ah         ; immediate 8 bit
ADD CX, 100
MOV AL, 'A'         ; AL <- 41h
\end{kodeasm}

\vspace{0.15em}
{\footnotesize\color{obeGray}Ukuran nilai mengikuti register tujuan:
8 bit untuk AL, 16 bit untuk AX.\par}
\Sumber{Bab XII, Bagian 1}
\end{frame}

% ---------------------------------------------------------
\begin{frame}[fragile]{Mode Direct Memory}
\small
Operand berada di memori; alamat (offset) ditulis di dalam instruksi.
Assembler mengganti nama variabel dengan offset sebenarnya.

\vspace{0.2em}
\begin{kodeasm}[numbers=none]
MOV AX, nilai       ; isi memori DS:nilai
MOV AX, [1234h]     ; isi memori pada offset 1234h
MOV total, AX       ; simpan AX ke variabel total
\end{kodeasm}

\vspace{0.15em}
{\footnotesize\color{obeGray}Cocok untuk variabel tunggal dengan
alamat tetap.\par}
\Sumber{Bab XII, Bagian 1}
\end{frame}

% ---------------------------------------------------------
\begin{frame}{Penulisan Mirip, Arti Berbeda}
\footnotesize
\begin{center}
\begin{tabular}{@{}>{\raggedright\arraybackslash}p{4.4cm}
    >{\raggedright\arraybackslash}p{8.6cm}@{}}
    \toprule
    \textbf{Penulisan} & \textbf{Arti} \\
    \midrule
    \code{MOV AX, nilai} & Menyalin isi variabel \code{nilai}
    (direct) \\
    \addlinespace
    \code{MOV AX, OFFSET nilai} & Menyalin offset variabel sebagai
    nilai segera \\
    \addlinespace
    \code{MOV AX, 1234h} & Menyalin nilai \hex{1234} \\
    \addlinespace
    \code{MOV AX, [1234h]} & Menyalin isi memori pada offset
    \hex{1234} \\
    \bottomrule
\end{tabular}
\end{center}

\vspace{0.2em}
\begin{alertblock}{Perhatikan tanda kurung}
\small Kurung siku mengubah arti dari nilai menjadi isi memori.
\end{alertblock}
\Sumber{Bab XII, Bagian 1}
\end{frame}

% ---------------------------------------------------------
\begin{frame}{Rangkuman Bagian 1}
\small
\begin{itemize}
    \item Mode pengalamatan menyatakan letak operand.

    \vspace{0.15em}
    \item Register: operand di register; immediate: nilai di
          instruksi.

    \vspace{0.15em}
    \item Direct memory: alamat tetap ditulis di instruksi.

    \vspace{0.15em}
    \item Tujuan tidak boleh berupa nilai segera.

    \vspace{0.15em}
    \item Kurung siku membedakan nilai dan isi memori.
\end{itemize}
\end{frame}

% =========================================================
% BAGIAN 2
% =========================================================
\PemisahBagian{Mode Memori Berpenunjuk}

% ---------------------------------------------------------
\begin{frame}[fragile]{Indirect Memory}
\small
Operand ditulis \code{[register]}; register menyimpan offset, bukan
data. Hanya BX, BP, SI, dan DI yang boleh menjadi penunjuk.

\vspace{0.2em}
\begin{kodeasm}[numbers=none]
MOV BX, OFFSET data1
MOV AX, [BX]        ; DS:BX
ADD BX, 2
MOV AX, [BX]        ; elemen kedua
\end{kodeasm}

\vspace{0.15em}
{\footnotesize\color{obeGray}BX/BP disebut basis; SI/DI disebut
indeks.\par}
\Sumber{Bab XII, Bagian 2}
\end{frame}

% ---------------------------------------------------------
\begin{frame}[fragile]{Base Relative dan Direct Index}
\small
\begin{itemize}
    \item \code{[BX+disp]} dan \code{[BP+disp]}: base relative.
    \item \code{[SI+disp]} dan \code{[DI+disp]}: direct index.
\end{itemize}

\vspace{0.2em}
\begin{kodeasm}[numbers=none]
MOV BX, OFFSET data1
MOV AX, [BX+4]      ; base relative
MOV SI, 4
MOV AX, data1[SI]   ; direct index
MOV AX, [SI+8]      ; direct index
\end{kodeasm}

\vspace{0.1em}
{\footnotesize\color{obeGray}\code{data1[SI]} diterjemahkan menjadi
\code{[data1+SI]}.\par}
\Sumber{Bab XII, Bagian 2}
\end{frame}

% ---------------------------------------------------------
\begin{frame}[fragile]{Based Indexed}
\small
Menggabungkan satu basis, satu indeks, dan satu perpindahan. Hanya
empat kombinasi yang legal:
\code{[BX+SI]}, \code{[BX+DI]}, \code{[BP+SI]}, \code{[BP+DI]}.

\vspace{0.2em}
\begin{kodeasm}[numbers=none]
MOV BX, OFFSET tabel
MOV SI, 6
MOV AX, [BX+SI+2]   ; medan ke-2 record ke-3
\end{kodeasm}

\vspace{0.1em}
{\footnotesize\color{obeGray}Cocok untuk matriks dan array of
record.\par}
\Sumber{Bab XII, Bagian 2}
\end{frame}

% ---------------------------------------------------------
\begin{frame}{Alamat Efektif}
\begin{center}
\begin{tikzpicture}[
    komp/.style={
        draw=obeBlue, line width=0.8pt, fill=obeBlueLight,
        rounded corners=3pt, minimum width=3.2cm,
        minimum height=0.8cm, align=center, font=\small
    },
    jumlah/.style={
        draw=obeGray, line width=0.9pt, fill=obeGrayLight,
        circle, minimum size=0.8cm, font=\large\bfseries
    },
    hasil/.style={
        draw=obeTeal, line width=0.9pt, fill=obeTealLight,
        rounded corners=3pt, minimum width=3.6cm,
        minimum height=0.9cm, align=center, font=\small
    },
    panah/.style={-{Stealth[length=2.2mm]}, line width=0.8pt, obeBlue}
]
    \node[komp] (base) at (0,2.2) {Basis};
    \node[komp] (index) at (0,1.2) {Indeks};
    \node[komp] (disp) at (0,0.2) {Perpindahan};
    \node[jumlah] (plus) at (4.0,1.2) {$+$};
    \node[hasil] (ea) at (7.6,1.2)
        {EA = basis + indeks + disp};
    \node[hasil] (fisik) at (7.6,-0.2)
        {Fisik = (segmen $\times$ 16) + EA};

    \draw[panah] (base) -- (plus);
    \draw[panah] (index) -- (plus);
    \draw[panah] (disp) -- (plus);
    \draw[panah] (plus) -- (ea);
    \draw[panah] (ea) -- (fisik);
\end{tikzpicture}
\end{center}

\vspace{0.2em}
{\footnotesize\color{obeGray}EA berupa offset 16 bit; alamat fisik
menambahkan basis segmen.\par}
\Sumber{Bab XII, Bagian 2}
\end{frame}

% ---------------------------------------------------------
\begin{frame}{Contoh EA dan Alamat Fisik}
\small
Diketahui \reg{DS} = \hex{1000}, \reg{BX} = \hex{0200},
\reg{SI} = \hex{0006}, dan \code{MOV AX, [BX+SI+2]}.

\[
\text{EA} = \hex{0200} + \hex{0006} + \hex{0002} = \hex{0208}
\]

\[
\text{alamat fisik} = (\hex{1000} \times 16) + \hex{0208}
= \hex{10000} + \hex{0208} = \hex{10208}
\]

\begin{block}{Inti}
\small EA hanya offset; alamat fisik memerlukan segmen.
\end{block}
\Sumber{Bab XII, Bagian 2}
\end{frame}

% ---------------------------------------------------------
\begin{frame}{Kombinasi Register Legal}
\scriptsize
\begin{center}
\begin{tabular}{@{}>{\raggedright\arraybackslash}p{3.2cm}
    >{\raggedright\arraybackslash}p{1.2cm}
    >{\raggedright\arraybackslash}p{1.4cm}
    >{\raggedright\arraybackslash}p{6.6cm}@{}}
    \toprule
    \textbf{Bentuk operand} & \textbf{Legal} & \textbf{Segmen} &
    \textbf{Alasan} \\
    \midrule
    \code{[BX]}, \code{[SI]}, \code{[DI]} & Ya & DS &
    Satu penunjuk \\
    \addlinespace
    \code{[BP]} & Ya & SS & BP adalah basis stack \\
    \addlinespace
    \code{[BX+SI]}, \code{[BX+DI]} & Ya & DS &
    Satu basis + satu indeks \\
    \addlinespace
    \code{[BP+SI]}, \code{[BP+DI]} & Ya & SS &
    Basis BP + satu indeks \\
    \addlinespace
    \code{[AX]}, \code{[CX]}, \code{[DX]}, \code{[SP]} & Tidak & --- &
    Bukan basis/indeks \\
    \addlinespace
    \code{[BX+BP]}, \code{[SI+DI]}, \code{[BX+SI+DI]} & Tidak & --- &
    Kombinasi tidak legal \\
    \bottomrule
\end{tabular}
\end{center}
\Sumber{Bab XII, Bagian 2}
\end{frame}

% ---------------------------------------------------------
\begin{frame}{Rangkuman Bagian 2}
\small
\begin{itemize}
    \item Indirect memakai satu register penunjuk: BX, BP, SI, atau
          DI.

    \vspace{0.15em}
    \item Base relative: basis + disp; direct index: indeks + disp.

    \vspace{0.15em}
    \item Based indexed: satu basis + satu indeks + disp.

    \vspace{0.15em}
    \item EA = basis + indeks + disp; alamat fisik menambahkan
          segmen.

    \vspace{0.15em}
    \item Hanya empat kombinasi based indexed yang legal.
\end{itemize}
\end{frame}

% =========================================================
% BAGIAN 3
% =========================================================
\PemisahBagian{Segmen Default dan Penerapan}

% ---------------------------------------------------------
\begin{frame}{Segmen Default}
\footnotesize
\begin{center}
\begin{tabular}{@{}>{\raggedright\arraybackslash}p{5.2cm}
    >{\raggedright\arraybackslash}p{1.8cm}
    >{\raggedright\arraybackslash}p{5.8cm}@{}}
    \toprule
    \textbf{Bentuk operand} & \textbf{Segmen} & \textbf{Keterangan} \\
    \midrule
    \code{[disp]}, nama variabel & DS & Data umum \\
    \addlinespace
    \code{[BX]}, \code{[SI]}, \code{[DI]}, \code{[BX+SI]},
    \code{[BX+disp]} & DS & Penunjuk data \\
    \addlinespace
    \code{[BP]}, \code{[BP+disp]}, \code{[BP+SI]},
    \code{[BP+DI]} & SS & Penunjuk stack \\
    \addlinespace
    Pengambilan instruksi (IP) & CS & Selalu segmen kode \\
    \bottomrule
\end{tabular}
\end{center}
\Sumber{Bab XII, Bagian 3}
\end{frame}

% ---------------------------------------------------------
\begin{frame}{Mengganti Segmen Default}
\small
\begin{itemize}
    \item Tuliskan register segmen di depan operand, misalnya
          \code{MOV AX, ES:[BX]}.

    \vspace{0.18em}
    \item Penggantian hanya mengubah register segmen; offset tetap.

    \vspace{0.18em}
    \item \code{[BP]} memakai SS, bukan DS.
\end{itemize}

\vspace{0.25em}
\begin{alertblock}{Jebakan \code{[BP]}}
\small \code{MOV AX, [BP]} membaca \eng{stack}. Untuk data biasa,
gunakan \code{DS:[BP]} atau register BX.
\end{alertblock}
\Sumber{Bab XII, Bagian 3}
\end{frame}

% ---------------------------------------------------------
\begin{frame}{Memilih Mode untuk Array dan Record}
\footnotesize
\begin{center}
\begin{tabular}{@{}>{\raggedright\arraybackslash}p{3.4cm}
    >{\raggedright\arraybackslash}p{4.4cm}
    >{\raggedright\arraybackslash}p{5.0cm}@{}}
    \toprule
    \textbf{Bentuk data} & \textbf{Mode yang tepat} &
    \textbf{Contoh} \\
    \midrule
    Variabel tunggal & Direct memory & \code{MOV AX, nilai} \\
    \addlinespace
    Array satu dimensi & Indirect atau direct index &
    \code{MOV AX, data1[SI]} \\
    \addlinespace
    Medan record & Base relative & \code{MOV AX, [BX+4]} \\
    \addlinespace
    Array of record/matriks & Based indexed &
    \code{MOV AX, [BX+SI+4]} \\
    \bottomrule
\end{tabular}
\end{center}
\Sumber{Bab XII, Bagian 3}
\end{frame}

% ---------------------------------------------------------
\begin{frame}{Rangkuman Bagian 3}
\small
\begin{itemize}
    \item Segmen default: DS untuk BX, SI, dan DI; SS untuk BP.

    \vspace{0.15em}
    \item Pengambilan instruksi selalu melalui CS:IP.

    \vspace{0.15em}
    \item Segmen dapat diganti dengan penulisan seperti
          \code{ES:[BX]}.

    \vspace{0.15em}
    \item Direct untuk variabel tunggal; indirect/direct index untuk
          array.

    \vspace{0.15em}
    \item Base relative untuk medan record; based indexed untuk
          array of record.
\end{itemize}
\end{frame}

% =========================================================
% CONTOH PROGRAM
% =========================================================
\begin{frame}[fragile]{Contoh 1: Tiga Mode untuk Array}
\begin{columns}[t]
\begin{column}{0.53\textwidth}
\begin{kodeasm}
.DATA
    data1 DW 10,20,30,40,50
    n     EQU 5
.CODE
START:
    MOV AX, @DATA
    MOV DS, AX
    MOV AX, data1+0
    ADD AX, data1+8
    MOV jumlah1, AX
    XOR SI, SI
    MOV CX, n
C2:
    ADD AX, data1[SI]
    ADD SI, 2
    LOOP C2
\end{kodeasm}
\end{column}
\begin{column}{0.44\textwidth}
\footnotesize
{\color{obeTeal}\bfseries Keluaran}\\[0.15em]
\begin{keluaran}
Perbandingan tiga mode pengalamatan
Ketiga hasil identik.
\end{keluaran}

\vspace{0.3em}
\begin{itemize}
    \item Direct: offset berupa konstanta.
    \item Direct index: \code{data1[SI]}, SI naik 2.
    \item Base relative: \code{[BX+SI]}.
    \item Jumlah kelima elemen = 150.
\end{itemize}
\end{column}
\end{columns}
\Sumber{Bab XII, Contoh 1}
\end{frame}

% ---------------------------------------------------------
\begin{frame}{Contoh 2: Tujuh Mode untuk Satu Elemen}
\scriptsize
\begin{center}
\begin{tabular}{@{}>{\raggedright\arraybackslash}p{2.4cm}
    >{\raggedright\arraybackslash}p{6.2cm}
    >{\raggedright\arraybackslash}p{4.0cm}@{}}
    \toprule
    \textbf{Mode} & \textbf{Potongan program} & \textbf{Alamat efektif} \\
    \midrule
    Register & \code{MOV AX, BX} & tanpa akses memori \\
    \addlinespace
    Immediate & \code{MOV AX, 30} & tanpa akses memori \\
    \addlinespace
    Direct & \code{MOV AX, data1+4} & offset data1 + 4 \\
    \addlinespace
    Indirect & \code{LEA SI, data1+4}; \code{MOV AX, [SI]} &
    offset data1 + 4 \\
    \addlinespace
    Base relative & \code{LEA BX, data1}; \code{MOV AX, [BX+4]} &
    offset data1 + 4 \\
    \addlinespace
    Direct index & \code{LEA SI, data1}; \code{MOV AX, [SI+4]} &
    offset data1 + 4 \\
    \addlinespace
    Based indexed & \code{LEA BX, data1}; \code{MOV SI, 4};
    \code{MOV AX, [BX+SI]} & offset data1 + 4 \\
    \bottomrule
\end{tabular}
\end{center}
\Sumber{Bab XII, Contoh 2}
\end{frame}

% =========================================================
% PRAKTIKUM
% =========================================================
\begin{frame}{Praktikum: Membandingkan Mode Pengalamatan}
\begin{block}{Judul praktikum}
\small Membandingkan mode pengalamatan pada akses array\\[0.15em]
\textbf{Sub-CPMK:} 3.3 \quad\textbullet\quad
\textbf{Waktu:} 100 menit \quad\textbullet\quad
\textbf{Bentuk:} individu
\end{block}

\vspace{0.2em}
\begin{columns}[t]
\begin{column}{0.46\textwidth}
{\color{obeBlue}\bfseries Alat dan bahan}
\begin{itemize}\footnotesize
    \item Komputer dengan emu8086
    \item DOSBox dengan TASM, TLINK, dan DEBUG
    \item Berkas \file{mode-array.asm}
    \item Lembar pengamatan register dan alamat
\end{itemize}
\end{column}
\begin{column}{0.50\textwidth}
{\color{obeBlue}\bfseries Langkah kerja}
\begin{enumerate}\footnotesize
    \item Rakit, link, dan jalankan program.
    \item Single step tiga mode; catat AX, SI, BX.
    \item Hitung EA setiap akses.
    \item Ubah isi array; pastikan hasil tetap sama.
    \item Periksa alamat array dengan DEBUG.
\end{enumerate}
\end{column}
\end{columns}

\vspace{0.2em}
{\footnotesize\textbf{Dikumpulkan:} berkas \file{.asm} modifikasi,
tangkapan layar, tabel EA, tangkapan DEBUG, dan penjelasan indirect
vs direct index.\par}
\Sumber{Bab XII, Praktikum}
\end{frame}

% =========================================================
% LATIHAN
% =========================================================
\begin{frame}{Latihan}
\small
\begin{enumerate}
    \item Jelaskan pengertian mode pengalamatan dan sebutkan tujuh
          mode yang dibahas.

    \vspace{0.15em}
    \item Tentukan mode untuk \code{MOV AX, BX},
          \code{MOV AX, 1234h}, \code{MOV AX, nilai},
          \code{MOV AX, [SI]}, \code{MOV AX, [BX+4]},
          \code{MOV AX, [DI+8]}, dan \code{MOV AX, [BX+SI+2]}.

    \vspace{0.15em}
    \item Jelaskan perbedaan \code{[BX+4]} dan \code{[SI+4]}.
          Mode apa yang dipakai masing-masing?

    \vspace{0.15em}
    \item Sebutkan empat kombinasi based indexed yang legal.
          Mengapa \code{[BX+BP]} dan \code{[SI+DI]} tidak legal?

    \vspace{0.15em}
    \item Diketahui DS = \hex{2000}, BX = \hex{0100}, SI =
          \hex{0006}. Hitung EA dan alamat fisik untuk
          \code{[BX+SI+4]}.
\end{enumerate}
\end{frame}

% =========================================================
% RANGKUMAN
% =========================================================
\begin{frame}{Rangkuman}
\small
\begin{itemize}
    \item Tujuh mode: register, immediate, direct, indirect, base
          relative, direct index, based indexed.

    \vspace{0.15em}
    \item Register dan immediate tidak mengakses memori.

    \vspace{0.15em}
    \item Indirect memakai satu register penunjuk.

    \vspace{0.15em}
    \item Base relative memakai basis + disp; direct index memakai
          indeks + disp.

    \vspace{0.15em}
    \item EA = basis + indeks + disp; alamat fisik menambahkan
          segmen.

    \vspace{0.15em}
    \item DS untuk BX/SI/DI; SS untuk BP.
\end{itemize}
\end{frame}

% ---------------------------------------------------------
\begin{frame}{Cek Pemahaman}
\footnotesize
\begin{enumerate}
    \item \code{MOV AX, [BX]} memakai mode \dots
          \begin{enumerate}
              \item register
              \item direct memory
              \item indirect memory
              \item immediate
          \end{enumerate}

    \vspace{0.2em}
    \item \code{MOV AX, [BX+SI+4]} memakai mode \dots
          \begin{enumerate}
              \item base relative \quad \item direct index \quad
              \item based indexed \quad \item indirect
          \end{enumerate}

    \vspace{0.2em}
    \item Segmen default untuk \code{MOV AX, [BP]} adalah \dots
          \begin{enumerate}
              \item CS \quad \item DS \quad
              \item SS \quad \item ES
          \end{enumerate}
\end{enumerate}
\end{frame}

% =========================================================
% PENUTUP
% =========================================================
\begin{frame}{Penutup}
\begin{block}{Sebelum pertemuan berikutnya}
\small
\begin{itemize}
    \item Selesaikan praktikum: berkas \file{.asm} modifikasi,
          tangkapan layar, tabel EA, dan penjelasan mode.
    \item Baca \textbf{Bab XIII} — Percabangan dan Perulangan.
    \item Siapkan pertanyaan tentang label, JMP, CMP, dan LOOP.
\end{itemize}
\end{block}

\vspace{0.4em}
\begin{center}
{\color{obeBlue}\sffamily\large\bfseries Terima kasih\par}
\vspace{0.15em}
{\color{obeGray}\sffamily\small \NamaMK\ (\KodeMK)
\textbullet\ \SemesterTeks\par}
\end{center}
\end{frame}

\end{document}
